\chapter[Sermon 100. Christ Is Shown Again to Be the Author of This Book, How Great He Is Here. Here Is Also Declared the Desire of the Church, Wishing for the Coming of Christ, and the Liberal Promise of the Lord.]{Christ Is Shown Again to Be the Author of This Book, How Great He Is Here. Here Is Also Declared the Desire of the Church, Wishing for the Coming of Christ, and the Liberal Promise of the Lord.}
\label{sermon:100}
\markright{Sermon 100. Christ Is Shown Again to Be the Author of This Book ...}

I Jesus sent my angel to testify to you these things in the congregations. I am the root and generation of David, and the bright morning star. And the Spirit and the bride said, “Come.” And let him who hears say also, “Come.” And let him who is thirsty come. And let whoever wills take of the water of life, freely.

\index{John, Saint}\index{Arethas of Caesarea}\index{Hebrews, Epistle to the}\index{Zurich}The tenth point of this conclusion again demonstrates that the author of this work is Jesus Christ, as Saint John states, to ensure that what is spoken carries greater authority and is more readily accepted by the audience. Therefore, nothing should be attributed to Saint John except for the writing of the work, that is to say, that he first saw these things and committed them to writing. The manner of the revelation is also repeated. Christ himself did not come down to earth, but sent forth his angel, who, from Christ and in Christ’s name, opened and showed these things to Saint John. The purpose of the angel’s sending or revelation is also specified: that he should testify these things in congregations and to all people throughout the world, until the end of the world. From these few words, we learn that credit must be given to this book, as it was proposed by the very Son of God through his angel and apostle, and indeed proposed to all within the church. Again, this demonstrates that Jesus Christ is very God, the Lord of angels, as Saint Paul affirms in the first chapter of Hebrews. This has also been spoken of before. These clear testimonies of Scripture ought to move the faithful more than the foolishness of Servetus the Spaniard and the Arian and Jewish practices of Servetans. Let us observe moreover that Christ sent his angel, not to judge or to teach, but to testify. Testimonies lawfully taken or committed to writing and sealed are not lawful to contradict; for they are entirely considered authentic. But this entire book was written by Saint John and is a witness or the testimony of God’s angel. Therefore, it is unlawful to doubt anything within it, and we ought to hold the same opinion of all other books of the Old and New Testaments. For the prophets and apostles are called the witnesses of God, and the gospel and doctrine of the prophets and apostles is the witness or testimony. He is mad who thinks that the canonical Scripture is of itself authentic unless it is first made authentic by the approval of the church and councils. Moreover, we understand that the doctrine of this entire book does not belong only to the seven churches of Asia, but to all who are scattered throughout the world; and therefore it pertains especially and singularly to us, who live today in Zurich or Switzerland, England, France, or Germany. Arethas, Bishop of Caesarea, says that he should testify, that is to say, that he should protest not privately or obscurely, but in the hearing of all churches dispersed throughout the world, so that no one pretending willful ignorance remains uncorrected.

\index{Ezekiel (Book of)}And immediately the Lord himself also shows and declares, who, and how great he is, and what we faithful have laid up in store in him. And he uses again parables and allusions for the more clarity. First, he calls himself the root and offspring of David, that is to say, a true and natural man. For we have heard before that he is very and natural God. And he cuts off from all heretics denying and impugning the true flesh of Christ, all witnesses. Strongly proving that he, after the flesh, is of our own nature. Whereof he is also called in Scripture the fruit of the womb of David, and he that is risen from his loins. Moreover, it is said to the Davidic virgin and mother of God: “You shall conceive in your womb and bring forth a son.” Therefore, he calls himself both the root and offspring of David. And the phrase of speech is to be marked. For the like is read in Ezekiel 16. “Your root and your offspring is of the land of Canaan”: that is to say, your birth is of the Canaanites, or your origin is of a polluted people. Yet, it seems here nevertheless also another certain thing is signified. For the root bears a tree and nourishes it, or quickens it. The root is not borne or nourished by the tree. And Christ the Lord is the foundation and preservation of the house of David and the church of the faithful. That David is preserved, that the offspring of David is not rooted out, which often times has deserved to be, it is done in respect or merit of Christ the Lord. Christ has saved them, the same saves also, so many as are saved, as he is the head of all the promises made unto David, adding virtue and even perfection, as in whom is perfect salvation and all fullness, as the clear testimonies of the Prophet Isaiah bear witness in chapters 7 and 37, and elsewhere, also in Hosea 3 and Ezekiel 34 and 37. And not an unlike place is in 1 Kings, chapter 15. John also in 1 [?John 1:1] of this book names Christ the root of David, and so forth.

\index{Apocalypse (Book of)}\index{Aquinas, Thomas}\index{Resurrection}Again the Lord calls himself a Star, and that not obscure, but shining and bright, even the morning Star. When he called himself a Star, he had respect to the most ancient oracle of Balaam, that most wise prophet in the East. He prophesied that a star should arise out of Israel, that is to say a celestial star, and even the very Son of God should be born of a woman. And that the same star did arise, the magicians testify in the second chapter of Saint Matthew. And it is called bright, because Christ is the light, illuminating all men that come into the world. Of this matter the same Saint John has treated much in the first, eighth, and ninth chapters of his Evangelical story. The same our Lord is also the morning star, so called by Saint Peter, 2 Peter 1, and by this our Saint John in the second chapter of the Apocalypse. For just as Lucifer rising, draws the day star after him, so Christ shining in the hearts of the faithful, does lighten them more and more in this present world also, and in the life to come does clothes them whole with the celestial light. Thomas of Aquin expounding this place says, “The morning star is to wit the messenger of the day, that is the everlasting felicity through his resurrection.” And these things we have heard hitherto from the mouth of Christ, concerning Christ, who and how great he is, and what treasures we have laid up in store in him. He is very God and man, was incarnate for us, that he might be our root, virtue, life, light, and salvation. Therefore have we reposed in him all fullness of salvation. And so we see again, that this book is written with the apostolic spirit, which spirit verily so often as occasion serves, reasons excellently of Christ, and preaches his salvation, and commends the faith in him, unto all the faithful. The same spirit therefore has inspired both book of the Gospel and the Apocalypse of Saint John, and caused them to be written by the same author.

11. In the eleventh place, the church is introduced, desiring the coming of Christ in judgment. For our Lord Jesus Christ is so good, so benign, and wholesome, that this entire book promises he will come and deliver the church of saints afflicted in this world. Now is recited the desire of this same church, wishing and calling the Lord, saying, “Come.” For soon we shall hear the Lord promising and saying, “I will come quickly.” And the church again responds, “Amen. Even so, come, Lord Jesus.” To show that the spirit within our body earnestly cries to the Lord for our deliverance and glorification, the Apostle mentions this much in Romans 8. By “the spirit,” however, we may understand every spiritual person also. And therefore Arethas, who calls them “spirit,” says that they are those accounted worthy of the spiritual marriage: and the bride, the church itself. Thus he says. We have spoken of the bride many times in this work, so we need not be tedious in repeating the same. Nevertheless, with a wonderful desire, all the godly long for the Lord to come in judgment: that day is terrible and abhorred for the wicked, but most joyful and wished for by the godly. For the godly perceive that they will once be delivered from all evils and abundantly rewarded with all good things, that the glory and truth of God will be advanced and established, that all ungodliness will be abolished, and the wicked will be tormented by the just judgment of God. Whereupon Peter in Acts 3 calls this day the restoration and fulfillment of all things that God has at any time spoken by the mouth of his prophets. In that same day, therefore, all the promises of God, even the greatest matters, will be fulfilled completely. Therefore, the Lord says in the gospel, “Lift up your heads, for your redemption draws near.” Those who mourn and are desperate cast down their heads; but the Lord bids us lift up our heads, to be cheerful and of good hope. For we shall certainly be delivered and glorified, who have been in the world a laughingstock and had in derision by all men. Therefore, the passages must be interpreted figuratively, which portray the exceedingly great lamentation and howling that will be in that day. For the wicked, in anguish, pain, and utter desperation, will cry out and tear themselves. The godly will rejoice in him whom they see coming, showing the wounds by which they are redeemed. Just as the desire of the saints was greatest when the first coming of our savior approached, as is evident in Simeon alone (Luke 2), so at the second coming of Christ in judgment, all the saints with unceasing voices will cry and continually do cry, “Come, Lord Jesus, come and deliver us, come and maintain your glory and church, almost brought to naught; come, our redeemer and savior, so wished and looked for, dispatch us from evils, grant us the promised good things, and so forth.”

Therefore, the following things may be referred either to the church or to John, that either the church or John should say: “Let him who hears say, ‘Come.’” Arethas explains this passage briefly and well: by these words, he suggests those who are not yet admitted to the flock, yet are ready to hear godly matters and diligently seek to know the Lord. So much he says. And without doubt, the desire of the godly is so great that they long for all creation to pray the Lord to come in judgment, as we often see in the Psalms, where the godly exhort the sun and moon and all creation to praise and speak well of the Lord.

\index{Antichrist}12. The twelfth place of the conclusion contains a most large promise and comfort of Christ. For he promises again plainly. As though he should say: I know what things the faithful will suffer under Antichrist, what also and how great craft the same will practice. All things he will sell for money, Heaven and Earth, and those things also which are not in his power. And he will deceive many, and will spoil many. And all the godly he will vex and oppress with grievous persecution. Therefore if I tarry long, and come not speedily, in as much as the desires of the saints covet the same, you that love and believe in me, flee Antichrist, give not yourselves to be spoiled by him. Look for me, have recourse unto me. He that is thirsty, that is, he that desires an heavenly gift, or he that is in anguish or tormented with cares and sundry evils, let him come to me; to me I say, let him come. I shall fill him with good things, deliver from evil, and will comfort him, and strengthen him with my spirit, in all manner dangers, that he may patiently bear and overcome all evils. And he seems to have borrowed these wholesome and most full of consolation words from the doctrine of Isaiah, which is in the fifty-fifth chapter, and in the seventh chapter of John. Hereof are spoken certain things about the beginning of the twenty-first chapter. Where we read the Lord to have said: And to him that is thirsty will I give of the well of the water of life freely.

\index{Primasius}But where he says, “and he who wills,” he does not mean, as many mistakenly believe, that it rests with our will to be saved. For we know that the Apostle has said, “It is not by will, nor by effort, but by God’s mercy.” The Lord saves us of his own good will, yet he does not save the unwilling, but the willing. But he gives us the ability to will, according to the Apostle’s saying, “It is God who works in us both to will and to accomplish.” Primasius says: by no good gifts does one go before receiving the water of life freely. For what do you have, says the Apostle, that you did not receive? Therefore, we have freely received from God the will to come also; to whom we gave nothing first, that we should be, much less that we should be made righteous from sinners. Thus he says. Nevertheless, it might seem to be a manner of speaking common among the Germans: which is, “I make it free for all to come,” I clearly exclude no one, I bid all come; so, “he who wills,” that is to say, “come all, and receive water,” and so forth. To the Lord be glory.
